% =============================================================================
% I. Introduction — IEEE-style skeleton (P2 draft)
% =============================================================================

\section{Introduction}
\label{sec:intro}

Multi-UAV cooperative navigation under partial observability and limited
communication is a core challenge for intelligent transportation and
urban air mobility.
Centralized training with decentralized execution (CTDE) methods such as MAPPO
provide stable local control, yet struggle to encode dynamic swarm interaction
structure when agents must coordinate through an evolving communication graph.

Graph neural networks and attention mechanisms have been widely adopted to
inject relational inductive bias into multi-agent reinforcement learning
(MARL).
In particular, GAT-style guidance can improve early coordination.
However, our empirical study shows a critical failure mode:
under prolonged training, graph-guided policies may \emph{increase training
reward while decreasing evaluation success}, indicating over-specialization to
surrogate alignment objectives rather than task-level generalization.
Dense Transformers further raise $O(N^{2})$ communication and compute costs,
which are poorly matched to range-limited airborne networks.

We therefore ask:
\begin{enumerate}
  \item Can spatiotemporal guidance improve long-horizon coordination without
        reward--success divergence?
  \item Can quality-aware sparse graphs retain performance under restricted
        communication budgets while controlling complexity?
\end{enumerate}

\paragraph{Contributions.}
This paper presents Dynamic Sparse Graph Fusion (DSGF) with the following
contributions:
\begin{itemize}
  \item A quality-weighted dynamic graph and sparse attention module for
        communication-aware UAV guidance, with per-step complexity
        $O(Nkd+Nd)$ when average degree $k\ll N$.
  \item A temporal residual policy interface that decouples guidance from
        surrogate reward shaping, mitigating long-horizon success collapse.
  \item A systematic IEEE-oriented evaluation covering main comparison,
        ablation, obstacle generalization, communication Pareto analysis, and
        scalability up to $32$ UAVs.
\end{itemize}

% Claim safety notes for writing:
% - Prefer "competitive under restricted communication" over "always lowest cost".
% - Never mix 4-UAV ablation percentages with 16-UAV Table I numbers.
% - Absolute success at N=16 is low for ALL methods; emphasize relative gains.
